\RequirePackage{pdfmanagement-testphase}
\DeclareDocumentMetadata{pdfversion=1.7}
\documentclass[11pt]{amsart}

\usepackage[margin=1in]{geometry}
\usepackage{amsmath,amssymb,amsthm,mathtools}
\usepackage[hidelinks]{hyperref}
\usepackage{microtype}
\usepackage{orcidlink}
\usepackage{tagpdf}
\tagpdfsetup{activate-all}

\hypersetup{
  pdftitle={Finite Soluble Groups Need Not Admit Increasing Unrefinable
    Subgroup Chains},
  pdfauthor={Richie Sater},
  pdfsubject={Maximal-subgroup index spectra in finite soluble groups},
  pdfkeywords={finite groups, maximal chains, maximal-subgroup indices,
    soluble groups, semidirect products}
}

\newtheorem{theorem}{Theorem}
\newtheorem{proposition}[theorem]{Proposition}
\newtheorem{lemma}[theorem]{Lemma}
\newtheorem{remark}[theorem]{Remark}
\newcommand{\F}{\mathbb F}
\newcommand{\I}{\mathcal I}

\title[Increasing Unrefinable Subgroup Chains]{Finite Soluble Groups Need Not
Admit Increasing Unrefinable Subgroup Chains}
\author[Richie Sater]{Richie Sater\,\orcidlink{0009-0007-9051-8207}}
\address{Independent Researcher, United States}
\email{richiesater@gmail.com}
\date{August 31, 2026}
\subjclass[2020]{Primary 20D30; Secondary 20D10, 20F22}
\keywords{finite groups, maximal chains, maximal-subgroup indices, soluble
groups, semidirect products}

\begin{document}
\begin{abstract}
An unrefinable subgroup chain is called increasing when its successive indices
are nondecreasing.  We show that not every finite group has such a chain, even
among soluble groups.  For every odd prime $p$, we construct a soluble matrix
group $G_p\leq \operatorname{GL}_5(p)$ of order $2^6p^8$ with no increasing
unrefinable chain.  The mechanism is a nested maximal-index-spectrum trap:
every candidate chain is forced through two steps of index $p^2$ to a subgroup
$N$, where the next admissible index is $2$.  All remaining lower indices
would then be $2$, contradicting $p\mid\lvert N\rvert$.  The construction
adapts module layers from an example of Kohler.  An ordinary group-theoretic
argument proves the result for all odd $p$; GAP calculations at $p=3,5,7$
corroborate selected instances.
\end{abstract}
\maketitle

\section{Introduction}

Let $G$ be a nontrivial finite group.  An unrefinable chain of subgroups
\[
  1=G_0<G_1<\cdots<G_n=G
\]
is one in which $G_{i-1}$ is maximal in $G_i$ for every $i$.  Associate with
this chain the indices $j_i=\lvert G_i:G_{i-1}\rvert$.  We call the chain
\emph{increasing} if
\[
  j_1\leq j_2\leq\cdots\leq j_n.
\]
Monakhov and Sokhor call such a chain a $<$-chain~\cite{MS2025}.  More
precisely, they prove that a finite group is supersoluble if it has two
maximal chains of the same length, one a $<$-chain and the other a
$>$-chain, and they classify the finite groups for which every maximal chain
in every proper subgroup is a $<$-chain or a $>$-chain.  In correspondence
beginning on August 8, 2026, they communicated to the author the following
question, stated here in equivalent terminology.\footnote{The correspondence is the source
through which the author learned of the question; no claim is made here about
its original formulation.  Reference~\cite{MS2025} supplies terminology and
background rather than a published source of this universal-existence
question.}
\begin{quote}
Does every finite group contain an increasing unrefinable subgroup chain?
\end{quote}

Kuperberg and Zieve study a different Jordan-type property for subgroup-lattice
intervals: under additional hypotheses, all maximal chains between a point
stabilizer and a transitive permutation group have the same relative indices
up to permutation~\cite{KuperbergZieve2007}.  Their invariance question across
all maximal chains is distinct from the present existential ordering question,
which asks whether some full chain has nondecreasing indices.

The condition is existential, so exhibiting a few decreasing chains cannot
answer the question.  A counterexample must control every possible maximal
step without requiring a description of the entire subgroup lattice.  We do
this by controlling the maximal-index spectra of three nested levels.

\begin{theorem}\label{thm:main}
For every odd prime $p$, there is a soluble group
$G_p\leq\operatorname{GL}_5(p)$ of order $2^6p^8$ which has no increasing
unrefinable subgroup chain.
\end{theorem}

The proof is organized around the following downward pattern, where the arrows
are maximal-subgroup steps and their labels are the corresponding indices:
\begin{equation}\label{eq:intro-trap}
 \begin{array}{ccccc}
 G_p & \xrightarrow{\ p^2\ } & M_X\text{ or }M_Y
     & \xrightarrow{\ p^2\ } & N,\\[-2pt]
 \I(G_p)=\{2,p^2\} &&
 \I(M_X)=\I(M_Y)=\{2,p^2,p^4\} &&
 \I(N)=\{2,p^4\}.
 \end{array}
\end{equation}
Reading downward, monotonicity forces a candidate chain through the two
index-$p^2$ steps to $N$.  There the only admissible next index is $2$, so
every remaining lower index would also be $2$.  This would make $N$ a
$2$-group, contradicting $p\mid\lvert N\rvert$.  Proposition
\ref{prop:trap} isolates this argument from the group construction.

The groups themselves adapt the $t=2$ case of Kohler's construction in
Section~4, especially Theorem~4.1, of~\cite{Kohler1964}.  Kohler used its
module structure to show that restrictions
on maximal-subgroup indices need not pass to subgroups.  Here two
nonisomorphic irreducible layers of dimension $2$ and an irreducible
commutator layer of dimension $4$ are used for a different purpose: they
realize the three spectra in \eqref{eq:intro-trap}.  Thus the reusable feature
is not the matrix presentation by itself, but the way successive module
quotients engineer a prescribed path of maximal indices.

The smallest displayed member is $G_3$, of order $2^6 3^8=419904$; no global
minimality is claimed.  The proof of Theorem~\ref{thm:main} is entirely
ordinary mathematics and uses no classification.  The repository's exact GAP
computations check the predicted spectra for $p=3,5,7$ and exhaustively search
the case $p=3$; they are finite corroboration rather than a proof of the
universal statement.  Section~2 gives the abstract obstruction, Section~3
constructs the required module layers, and Section~4 computes their spectra
and proves the theorem.  Further questions are recorded in Section~5, while
the explicit matrices and computational boundary appear in Appendix~A.

\section{A nested index-spectrum trap}

For a finite group $A$, write
\[
  \I(A)=\{\lvert A:B\rvert\mid B\text{ is maximal in }A\}.
\]
The next proposition is the conceptual core of the proof.  It replaces a
global search through the subgroup lattice by three local spectra and two
incidence statements.

\begin{proposition}[Nested spectrum trap]\label{prop:trap}
Let $p$ be an odd prime, and let $G$ be a finite group with subgroups
$M_1,M_2,N$ such that $M_1$ and $M_2$ are maximal in $G$, while $N$ is
maximal in each $M_i$.  Suppose that
\begin{align*}
 \I(G)&=\{2,p^2\},\\
 \I(M_1)=\I(M_2)&=\{2,p^2,p^4\},\\
 \I(N)&=\{2,p^4\}.
\end{align*}
Assume also that
\begin{enumerate}
\item every maximal subgroup of $G$ of index $p^2$ is $G$-conjugate to
      $M_1$ or $M_2$;
\item for $i=1,2$, every maximal subgroup of $M_i$ of index $p^2$ is
      $M_i$-conjugate to $N$; and
\item $p$ divides $\lvert N\rvert$.
\end{enumerate}
Then $G$ has no increasing unrefinable subgroup chain.
\end{proposition}

\begin{proof}
Suppose that
\[
  1=G_0<G_1<\cdots<G_n=G
\]
is unrefinable and that $j_i=\lvert G_i:G_{i-1}\rvert$ is nondecreasing.
The last index belongs to $\I(G)$.  If $j_n=2$, then every $j_i$ equals $2$,
so $\lvert G\rvert$ is a power of $2$; this contradicts
$p\mid\lvert N\rvert$ and $N\leq G$.  Hence $j_n=p^2$.  After conjugating
the chain in $G$, we may take $G_{n-1}=M_i$ for some $i\in\{1,2\}$.

Now $j_{n-1}\leq p^2$.  The spectrum of $M_i$ leaves only $2$ and $p^2$.
The value $2$ would again force every preceding index to be $2$, contrary to
$p\mid\lvert N\rvert$ and $N\leq M_i$.  Thus $j_{n-1}=p^2$.  Conjugating
the lower part of the chain inside $M_i$, we may take $G_{n-2}=N$.

Finally, $j_{n-2}\leq p^2$, whereas $\I(N)=\{2,p^4\}$.  Therefore
$j_{n-2}=2$, and monotonicity forces all earlier indices to equal $2$.
This would make $\lvert N\rvert$ a power of $2$, contradicting
$p\mid\lvert N\rvert$.
\end{proof}

The proof uses only the labeled local data
\[
  \{2,p^2\}\longrightarrow\{2,p^2,p^4\}
  \longrightarrow\{2,p^4\}.
\]
In particular, one need not enumerate maximal chains or even determine the
full subgroup lattice.  The remaining task is to realize this diagram in a
soluble group.

\section{The group family and its module layers}

We first record the standard maximal-subgroup fact that converts irreducible
module dimensions into subgroup indices.  We then construct three layers of
dimensions $2,2,4$, exactly the dimensions required by Proposition
\ref{prop:trap}.

\subsection{Maximal subgroups of a coprime semidirect product}

\begin{lemma}\label{lem:semidirect}
Let $V$ be an elementary abelian $p$-group, let $H$ be a $p'$-group acting
on $V$, put $K=V\rtimes H$, and let $\pi:K\to H$ be the natural projection.
Every maximal subgroup of $K$ has one of the following forms:
\begin{enumerate}
\item $\pi^{-1}(J)=V\rtimes J$, where $J$ is maximal in $H$;
\item a $V$-conjugate of $W\rtimes H$, where $W$ is a maximal proper
      $H$-submodule of $V$.
\end{enumerate}
\end{lemma}

\begin{proof}
Let $T$ be maximal in $K$.  If $V\leq T$, then $T/V$ is maximal in
$K/V\cong H$, which gives the first form.  Suppose $V\nleq T$.  Then
$VT=K$.  Put $W=V\cap T$.  The group $T$ normalizes $W$, while $V$
centralizes $W$; hence $W\lhd K$.  In $K/W$, both $T/W$ and $WH/W$ are
complements to $V/W$.  By the Schur--Zassenhaus theorem, these complements
are conjugate by an element of $V/W$.  Lifting this conjugacy shows that $T$
is $V$-conjugate to $W\rtimes H$.  Moreover, maximality of $T$ is equivalent
to maximality of $W$ among the proper $H$-submodules of $V$.
\end{proof}

If $H$ is a $2$-group, the subgroups of the first type have index $2$.  If
$V$ is completely reducible, the indices of the second type are the orders
of the irreducible quotients of $V$.  This is the interface used below: an
irreducible quotient of dimension $d$ over $\F_p$ produces index $p^d$.

\subsection{The matrix group}

Fix an odd prime $p$ and put $F=\F_p$.  Let
\[
 s=\begin{pmatrix}0&1\\1&0\end{pmatrix},\qquad
 t=\begin{pmatrix}1&0\\0&-1\end{pmatrix},\qquad
 D=\langle s,t\rangle\leq\operatorname{GL}_2(F).
\]
Then $D\cong D_8$, where $D_8$ has order $8$.  The natural $FD$-module
$U=F^2$ is absolutely irreducible.  Indeed, after extending scalars to an
algebraic closure, the only $t$-invariant lines are the two coordinate lines,
and $s$ interchanges them.

Let $H=D\times D$.  The $p$-group on which $H$ will act is
\[
 L=\left\{
 \ell(x,y,z)=
 \begin{pmatrix}
 I_2&x&z\\
 0&1&y\\
 0&0&I_2
 \end{pmatrix}\mathrel{:}
 x\in F^{2\times1},\ y\in F^{1\times2},\ z\in F^{2\times2}
 \right\}.
\]
Here the block sizes are $2,1,2$.  The following two calculations identify
the degree-one coordinates and their commutator layer.  Direct multiplication
gives
\begin{equation}
 \ell(x,y,z)\ell(x',y',z')
 =\ell(x+x',y+y',z+z'+xy'),
\end{equation}
and, with $[a,b]=a^{-1}b^{-1}ab$,
\begin{equation}\label{eq:commutator}
 [\ell(x,y,z),\ell(x',y',z')]
 =\ell(0,0,xy'-x'y).
\end{equation}
Thus $\lvert L\rvert=p^8$.  Put
\[
 Z=\{\ell(0,0,z)\},\qquad
 X=\{\ell(x,0,0)\},\qquad
 Y=\{\ell(0,y,0)\}.
\]
The group $Z$ is central in $L$.  The outer products $xy$ span
$F^{2\times2}$, so \eqref{eq:commutator} gives $L'=Z$.  Since $p$ is odd,
$L$ has exponent $p$; explicitly,
\[
 \ell(x,y,z)^p
 =\ell\left(0,0,pz+\binom p2xy\right)=1.
\]
The standard formula for the Frattini subgroup of a finite $p$-group now
gives
\begin{equation}\label{eq:frattini}
 \Phi(L)=L^pL'=Z.
\end{equation}
The four coordinates in $x$ and $y$ and the four independent coordinates in
$z$ also make the order $p^{4+4}=p^8$ transparent.

Embed $H$ in $\operatorname{GL}_5(F)$ by identifying $(A,B)\in D\times D$
with $\operatorname{diag}(A,1,B)$.  Conjugation acts by
\[
 x\mapsto Ax,\qquad y\mapsto yB^{-1},\qquad z\mapsto AzB^{-1}.
\]
Consequently, $H$ normalizes both $L$ and $Z$.  Define
\[
 G_p=L\rtimes H\leq\operatorname{GL}_5(p).
\]
Then $Z\lhd G_p$ and
$\lvert G_p\rvert=p^8\cdot 8^2=2^6p^8$.  Since both $L$ and $H$ are
soluble, so is $G_p$.

\subsection{The layers of dimensions \texorpdfstring{$2,2,4$}{2,2,4}}

Set $\overline X=XZ/Z$ and $\overline Y=YZ/Z$.  As an $H$-module,
\begin{equation}\label{eq:modules}
 L/Z=\overline X\oplus\overline Y,
\end{equation}
where $\overline X$ and $\overline Y$ are nonisomorphic irreducible modules
of dimension $2$.  The first factor of $D\times D$ acts naturally on
$\overline X$ and the second acts trivially; the second factor acts dually on
$\overline Y$ and the first acts trivially.  In particular,
\[
 C_H(\overline X)=1\times D,
 \qquad
 C_H(\overline Y)=D\times 1.
\]
The distinct action kernels show that the modules are nonisomorphic.  Since
$p\nmid\lvert H\rvert=64$, Maschke's theorem applies, and the only maximal
$H$-submodules of $L/Z$ are $\overline X$ and $\overline Y$.

The commutator layer is
\begin{equation}
 Z\cong U\boxtimes U^*,
\end{equation}
the external tensor product for the two direct factors of $H$.  Hence $Z$ is
an irreducible $H$-module of dimension $4$.  To see this over the given
field, extend scalars to an algebraic closure $\overline F$: both
$U\otimes_F\overline F$ and $U^*\otimes_F\overline F$ remain irreducible,
and their external tensor product is irreducible for $D\times D$.

The role of these dimensions can now be read backward from the desired trap.
The two dimension-$2$ quotients in $L/Z$ create the two top steps of index
$p^2$.  After one such quotient is removed, the remaining elementary abelian
normal subgroup has irreducible layers of dimensions $2$ and $4$.  Removing
the second dimension-$2$ layer leaves the irreducible module $Z$, whose only
non-$2$ maximal index is $p^4$.  This is precisely the transition displayed
in \eqref{eq:intro-trap}.

\section{Realizing the spectrum trap}

We now compute only the maximal-subgroup data required by Proposition
\ref{prop:trap}.  The calculation has three levels: the quotient $G_p/Z$,
the two middle semidirect products, and the terminal group $Z\rtimes H$.

\begin{lemma}\label{lem:spectra}
Put
\[
 M_X=(Z\times X)\rtimes H,\qquad
 M_Y=(Z\times Y)\rtimes H,\qquad
 N=Z\rtimes H.
\]
The subgroups $M_X$ and $M_Y$ are maximal in $G_p$ of index $p^2$, and $N$
is maximal in each of them with index $p^2$.  Moreover,
\begin{align*}
 \I(G_p)&=\{2,p^2\},\\
 \I(M_X)=\I(M_Y)&=\{2,p^2,p^4\},\\
 \I(N)&=\{2,p^4\}.
\end{align*}
Every maximal subgroup of $G_p$ of index $p^2$ is $G_p$-conjugate to $M_X$
or $M_Y$.  Every maximal subgroup of $M_X$ of index $p^2$ is
$M_X$-conjugate to $N$, and every maximal subgroup of $M_Y$ of index $p^2$
is $M_Y$-conjugate to $N$.
\end{lemma}

\begin{proof}
\emph{The top level.}
First, every maximal subgroup $M$ of $G_p$ contains $Z$.  If $Z\nleq M$,
then maximality gives $G_p=ZM$, and the modular law together with
\eqref{eq:frattini} gives
\[
 L=L\cap ZM=Z(L\cap M)=\Phi(L)(L\cap M).
\]
Recall that $L=\Phi(L)K$ for a subgroup $K\leq L$ implies $K=L$: otherwise,
a maximal subgroup of $L$ containing $K$ would contain both $K$ and
$\Phi(L)$, a contradiction.  Applying this fact to $K=L\cap M$ forces
$L\cap M=L$, so $Z\leq M$, again a contradiction.

It follows that the maximal subgroups of $G_p$ are the inverse images of the
maximal subgroups of
\[
 G_p/Z=(\overline X\oplus\overline Y)\rtimes H.
\]
Lemma~\ref{lem:semidirect} and \eqref{eq:modules} give index $2$ from maximal
subgroups of the $2$-group $H$, and index $p^2$ from either irreducible
quotient.  They also show that the index-$p^2$ maximal subgroups are conjugate
to $M_X$ or $M_Y$.  This proves the top spectrum and its incidence statement.

\emph{The middle level.}
In $M_X$, the normal elementary abelian subgroup $Z\times X$ is the direct
sum of nonisomorphic irreducible $H$-modules of dimensions $4$ and $2$.  Its
only maximal $H$-submodules are therefore $Z$ and $X$.  Lemma
\ref{lem:semidirect} gives the indices $2,p^2,p^4$; the index-$p^2$ case
retains $Z$ and hence gives an $M_X$-conjugate of $N$.  The same argument
applies to $M_Y$.

\emph{The terminal level.}
The $H$-module $Z$ is irreducible of dimension $4$.  Applying Lemma
\ref{lem:semidirect} to $N=Z\rtimes H$ gives
$\I(N)=\{2,p^4\}$.  These three calculations establish every spectrum and
incidence needed for the trap.
\end{proof}

\begin{proof}[Proof of Theorem~\ref{thm:main}]
The construction gives a soluble group $G_p$ of order $2^6p^8$.  Lemma
\ref{lem:spectra} supplies the three spectra and both incidence statements in
Proposition~\ref{prop:trap}; moreover, $p$ divides
$\lvert N\rvert=2^6p^4$.  The proposition therefore excludes every
increasing unrefinable subgroup chain in $G_p$.
\end{proof}

The proof has thus separated the construction from the obstruction.  Module
theory establishes the local diagram, and Proposition~\ref{prop:trap} turns
that diagram into the global nonexistence result.

\section{Conclusion}

The spectrum trap is a small directed system rather than a subgroup-lattice
classification: its vertices record maximal-index spectra, and its marked
arrows record which maximal subgroups can follow an admissible index.  This
viewpoint may be useful whenever one wants to force prescribed rises or drops
in unrefinable chains.  Different group constructions could realize the same
local diagram.

In this example, the two nonisomorphic dimension-$2$ layers and the
irreducible dimension-$4$ layer are structurally essential to the displayed
pattern.  Coprime action supplies complete reducibility, while the
$2$-group complement makes every quotient-side maximal index equal to $2$.
Neither a classification theorem nor a computation enters the uniform proof.

\medskip
\noindent\textbf{Further questions.}
Three natural questions remain.  What is the least order of a counterexample?
Which labeled spectrum diagrams can be realized by finite soluble groups?
And can similar module-layer constructions prescribe more complicated local
patterns in maximal chains?  The present argument neither establishes the
minimality of $G_3$ nor classifies groups that admit increasing chains.

\appendix
\section{The explicit group \texorpdfstring{$G_3$}{G3} and the computational boundary}

This appendix gives one concrete representative and states exactly what the
repository computations check.  It is not needed for the proof for arbitrary
odd $p$.

All matrices below are over $\F_3$, so $-1=2$.  Let $E_{ij}$ denote the
standard matrix unit, write $e_{ij}=I_5+E_{ij}$, and set
\[
 \begin{array}{ll}
 a_s=\operatorname{diag}(s,1,I_2),&
 a_t=\operatorname{diag}(t,1,I_2),\\
 b_s=\operatorname{diag}(I_2,1,s),&
 b_t=\operatorname{diag}(I_2,1,t).
 \end{array}
\]
Then the smallest displayed example is
\begin{equation}\label{eq:explicit-g3}
 G_3=\langle e_{13},e_{23},e_{34},e_{35},a_s,a_t,b_s,b_t\rangle
 \leq\operatorname{GL}_5(3).
\end{equation}
Define
\[
 \begin{split}
 Z_3&=\langle e_{14},e_{15},e_{24},e_{25}\rangle,\\
 H_3&=\langle a_s,a_t,b_s,b_t\rangle,\\
 M_{X,3}&=\langle Z_3,e_{13},e_{23},H_3\rangle,\\
 M_{Y,3}&=\langle Z_3,e_{34},e_{35},H_3\rangle,
 \qquad N_3=\langle Z_3,H_3\rangle.
 \end{split}
\]
Here $\lvert H_3\rvert=64$.  The relevant orders and complete
maximal-index spectra are
\[
\begin{array}{c|c|c}
 A&\lvert A\rvert&\I(A)\\ \hline
 G_3&419904&\{2,9\}\\
 M_{X,3},M_{Y,3}&46656&\{2,9,81\}\\
 N_3&5184&\{2,81\}.
\end{array}
\]

The ordinary argument in Sections~2--4 proves Theorem~\ref{thm:main} for
every odd prime.  The standalone file \texttt{G3-verification.g} separately
defines the eight generators in \eqref{eq:explicit-g3} as literal
$5\times5$ matrices.  Using standard GAP commands and no optional package, it
checks the named orders and maximal inclusions, enumerates every conjugacy
class of maximal subgroups in the four relevant groups, checks coverage of
the marked conjugacy classes, and performs an exhaustive top-down search for
an increasing chain in $G_3$.  Captured successful runs use GAP 4.11.1 and
GAP 4.16.0.

The family tests in \texttt{tests/counterexample.g} instantiate $p=3,5,7$
and reproduce the predicted spectra and incidences.  The separate
\texttt{tests/mmc-crosscheck.g} compares the repository's chain recurrence
with exhaustive index-sequence enumeration for all $144$ groups of order at
most $32$.  This last test checks the implementation of the recurrence, not
the theorem's general construction.  Together these exact computations can
catch transcription and implementation errors in the displayed instances;
they do not replace the uniform proof.

The theorem-related GAP scripts and captured outputs described above are
archived in release
\href{https://github.com/RichieSater/monotone-maximal-chains/releases/tag/v2.0.0}{v2.0.0}
of the accompanying repository~\cite{Repository}.

\section*{Acknowledgments}

The author thanks V.~S.~Monakhov and I.~L.~Sokhor for bringing the question
to his attention through their correspondence.

\section*{Code and data availability}

The computational supplement is publicly available in the
\href{https://github.com/RichieSater/monotone-maximal-chains}{verification
repository}.  Release \texttt{v2.0.0} contains the standalone verifier and
captured outputs from GAP 4.11.1 and GAP 4.16.0.  Source code is licensed
under the MIT License, and the manuscript and documentation are licensed
under CC BY 4.0.

\section*{Generative-AI disclosure}

From 8 to 31 August 2026, the author used the
\href{https://developers.openai.com/codex/cli/}{OpenAI Codex CLI} (versions
0.147.0, 0.150.1, and 0.151.0; model GPT-5.6-sol; accessed locally via the
command line) for mathematical exploration, literature searches, GAP and
regression-test development, proof auditing, and editorial revision; the
mathematical argument here proves the theorem, the reported computations have
their stated finite scope, and the author takes responsibility for the manuscript.

\begin{thebibliography}{9}

\bibitem{Kohler1964}
J. Kohler,
\emph{Finite groups with all maximal subgroups of prime or prime square
index},
Canad. J. Math. \textbf{16} (1964), 435--442.
\href{https://doi.org/10.4153/CJM-1964-046-6}{doi:10.4153/CJM-1964-046-6}.

\bibitem{KuperbergZieve2007}
G. Kuperberg and M. E. Zieve,
\emph{Analogues of the Jordan--H\"older theorem for transitive $G$-sets},
arXiv:0712.4142 (2007).
\href{https://doi.org/10.48550/arXiv.0712.4142}{doi:10.48550/arXiv.0712.4142}.

\bibitem{MS2025}
V. S. Monakhov and I. L. Sokhor,
\emph{On indices of maximal chains in finite groups},
Results Math. \textbf{80} (2025), article 155.
\href{https://doi.org/10.1007/s00025-025-02468-5}{doi:10.1007/s00025-025-02468-5}.

\bibitem{Repository}
R. Sater,
\emph{Finite soluble groups need not admit increasing unrefinable subgroup
chains}, version 2.0.0 preprint and verification archive (2026),
\href{https://doi.org/10.5281/zenodo.22213657}{doi:10.5281/zenodo.22213657}.

\end{thebibliography}

\end{document}
